\def \Nperiod {{N_0}} %period of discrete-time signals, avoid to use N
\def \Tperiod {{T_0}} %period of continuous-time signals, avoid to use T

%\def \digif {{\mathbb{F}}} %discrete-time linear frequency %AK: I stopped using it


%from http://tex.stackexchange.com/questions/42726/align-but-show-one-equation-number-at-the-end
\newcommand\numberthis{\addtocounter{equation}{1}\tag{\theequation}}

\def \nuv {{\boldsymbol{\nu}}} %boldface \nu. assumes amsmath package as in  http://tex.stackexchange.com/questions/595/how-can-i-get-bold-math-symbols

\def \cconv {{\circledast}} %\varoast \oast %circular convolution

\newcommand{\figwidth}{.7\textwidth} %defaut fig width
\newcommand{\figwidthSmall}{.5\textwidth} %smaller fig width
\newcommand{\figwidthLarge}{.9\textwidth} %larger fig width

\newcommand{{\ip}}[2]{{ \langle {#1} , {#2} \rangle }} %inner product. Usage example: $\ip{\varphi_j(t)}{\varphi_k(t)} = 1$

%This command allows to use \href when compiling with pdflatex or disable it when compiling for a DVI with latex
%I use it to create a list of URLs, which I will later identify on a web site via urls_akbook.tex
\newcommand{{\akurl}}[2]{{\ifpdf{[\href{{#1}}{{url#2}}]}\else{\IfPackageLoadedTF{tex4ht}{[\href{{#1}}{{url#2}}]}{[url#2]}}\fi}} %

\def \sign {{\textrm{sign}}} %sign function

\newcommand{\sgn}{\mathop{\mathrm{sgn}}}

\def \aktilde {{$\sim$}} 
%\def \aktilde {{\textrm{\textasciitilde}}} %this is too high, there is a nice \midtilde in wsuipa package but did not want to include the package
\def \muV {{$\mu$V}}  %micro volts (1e-6)
\def \mus {{$\mu$s}}  %micro seconds (1e-6)

%distinguish digital and analog frequencies
%\def \dw {{\omega}}  %in radians
%\def \aw {{\Omega}}  %in radians/sec
\def \aw {{\omega}}  %in radians
\def \dw {{\Omega}}  %in radians/sec

%############

\def \bModel     {\bthm{Model}}
\def \eModel     {\ethm{Model}}

%#######################
%Shortcuts
%#######################
\def \xrms {{x_{\textrm{rms}}}} %RMS value of x
\def \BW {{\textrm{BW}}} %bandwidth
\def \MSE {{\textrm{MSE}}} %mean-squared error
\def \SNR {{\textrm{SNR}}} %SNR

\def \dB {{\textrm{dB}}} %decibel
\def \sp {{\textrm{~}}}

%\newcommand {\arrowedbox}[1] {{\rightarrow\fbox{#1}\rightarrow}}
\def \setM {{\mathscr M}}

%\real and \imag for real and imaginary parts
\newcommand{\real}[1]{{\mathcal Real \{#1\}}}
\newcommand{\imag}[1]{{\mathcal Imag \{#1\}}}


%Similarly, discussions on comp.text.tex have revealed that there are a variety of
%ways to indicate the mathematical notion of "is defined as". Common candidates include:
%\triangleq), \equiv), \coloneqq) and:
%\def \define {{{\stackrel{\text{\tiny def}}{=}}}}  %def gets too close of left-side
%\def \define {{\triangleq}}

\def \cm {{Common mistake}}

\def \tmax {{{\textrm{max}}}}
\def \tmin {{{\textrm{min}}}}

%random signals
\def \rsx {{\mathsf x}}
\def \rsy {{\mathsf y}}
\def \rsz {{\mathsf z}}

%random variables
\def \rvn {{\mathsf N}}
\def \rvr {{\mathsf R}}
\def \rvk {{\mathsf K}}
\def \rve {{\mathsf E}}
\def \rvx {{\mathsf X}}
\def \rvy {{\mathsf Y}}
\def \rvz {{\mathsf Z}}
%\def \rvx {{\mathbb X}}
%\def \rvy {{\mathbb Y}}
%\def \rvz {{\mathbb Z}}

\def \conv {{\ast}}  %convolution
%\def \conv {{\convolution}}  %package mathabx convolution

%moved to definitions.tex
%use labels starting with eq: and fig: for equations and figures
%\newcommand{\blol}[1] {{Block~(\ref{eq:#1})}}
%\renewcommand{\equl}[1] {{Eq.~(\ref{eq:#1})}}
%\renewcommand{\figl}[1] {{Figure~\ref{fig:#1}}}
%\newcommand{\tabl}[1] {{Table~\ref{tab:#1}}}
%use labels starting with ak_ for listings
%\newcommand{\codl}[1] {{Listing~\ref{code:#1}}}
%\newcommand{\exal}[1] {{Example~\ref{ex:#1}}}

%expected value:
\renewcommand{\ev}{\mathbb{E}}
%\def \expval {\textbf{E}}
\def \expval {\ev}

%#######################
%Need to distinguish
%#######################
\def \ts {{T_s}}%sampling period
\def \fs {{F_s}}%sampling frequency
\def \tsym {{T_{\textrm{sym}}}}%symbol period
\def \rsym {{R_{\textrm{sym}}}} %symbol rate
\def \tdmt {{T_{\textrm{dmt}}}}%symbol period
\def \rdmt {{R_{\textrm{dmt}}}} %symbol rate
\def \xa {{$x(t)$}} %analog
\def \xd {{$x_q[n]$}}  %digital
\def \xs {{$x_s(t)$}} %sampled
\def \xsarea#1{{{x_s^{\textrm{area}}({{#1}})}}} %area of sampled signal

%#######################
%Code listing
%#######################
%Making easier to include the listing of source code:
%assume name starts with ak_
%Example:
%directory, name
%\includecode{MatlabOctaveFunctions}{dctmtx}
%is equivalent to
%\lstinputlisting[caption={MatlabOctaveFunctions/ak\_dctmtx.m},label=code:dctmtx]{./Code/MatlabOctaveFunctions/ak_dctmtx.m}
\def \includecode#1#2 {{\lstinputlisting[caption={#1/ak\_#2.m},label=code:#2]{./Code/#1/ak_#2.m}}}

%need to take care of the cases where an underscore _ exists or the cases where ak is not the prefix
%directory, caption, name
\def \includecodelong#1#2#3 {{\lstinputlisting[caption={#1/#2.m},label=code:#3]{./Code/#1/#3.m}}}
%Examples:
%\includecodelong{MatlabOctaveFunctions}{qpsk\_simple}{qpsk_simple}
%\includecodelong{MatlabOctaveFunctions}{ak\_gram\_schmidt}{ak_gram_schmidt}
%In case the macros do not work, use:
%\lstinputlisting[language=Python,caption=MatlabBookFigures/figs\_transforms\_dctimagecoding,label=code:dctimagecoding]{./Code/MatlabBookFigures/figs_transforms_dctimagecoding.m}

%now add hyperlink to Python version
\newcommand{{\akpythoncode}}[2]{{\ifpdf{[\href{{#1}}{{#2}}]}\else{\IfPackageLoadedTF{tex4ht}{[\href{{#1}}{{#2}}]}{[#1]}}\fi}} %

\def \includecodepython#1#2#3 {{\lstinputlisting[caption={#1/#2.m. \akpythoncode{https://github.com/aldebaro/dsp-telecom-book-code/blob/master/PythonCodeSnippets/#2.py}{Python version}},label=code:#3]{./Code/#1/#3.m} }}


%need to take care of the cases where an underscore _ exists or the cases where ak is not the prefix
%directory, caption, name
\def \onlypythoncodelong#1#2#3 {{\lstinputlisting[caption={#1/#2.py},label=code:#3]{./Code/#1/#3.py}}}


% #### AK: the label cannot have underscore: _
%\def \onlypythoncodelonglines#1#2#3#4#5#6 {{\lstinputlisting[caption={Method #6 in lasse-py/lasse/#1/#2.py},label=code:#6,language=Python,firstline=#4,lastline=#5]{../../lasse-py/lasse/#1/#3.py}}}
\newcommand{\onlypythoncodelonglines}[7]{%
\lstinputlisting[
  caption={Method \texttt{#6} in \texttt{lasse-py/lasse/#1/#2.py}},
  label={code:#7},
  language=Python,
  firstline=#4,
  lastline=#5
]{../../lasse-py/lasse/#1/#3.py}%
}


%First input argument is the folder and the second is the file name after ak_
%For instance:
%\onlypythoncode{PythonFunctions}{convolution}
%finds the file ./Code/PythonFunctions/ak_convolution.py
\newcommand{\onlypythoncode}[2]{%
  \lstinputlisting[
    caption={#1/ak\_#2.py},
    label={code:#2},
    language=Python
  ]{./Code/#1/ak_#2.py}%
}

\newcommand{\onlypythonwithlines}[4]{%
  \lstinputlisting[
    caption={#1/ak\_#2.py},
    label={code:#2},
    language=Python,
		firstline=#3,
		lastline=#4
  ]{./Code/#1/ak_#2.py}%
}


%#######################
%Others
%#######################

%three-dimensional array "channel matrix"
\newcommand{\cmH}{{ \overline {\textbf H}}}

\def \matlab {{{Matlab/Octave}}}

%total and max
%\def \ptot {{p^{\textrm{tot}}}}
\def \ptot {{{\calP}^{\textrm{tot}}}}

\def \pmax {{P^{\textrm{max}}}}
\def \rtarget {{r^{\textrm{target}}}}
%allow to provide subscript
\newcommand{\ptotn}[1] {{p_{#1}^{\textrm{tot}}}}
\newcommand{\pmaxn}[1] {{P_{#1}^{\textrm{max}}}}
\newcommand{\rtargetn}[1] {{r_{#1}^{\textrm{tar}}}}

\def \grandmatrix {{\overline{\mathbf G}}}

\newcommand{\nk}[1] {{#1_n^k}}

\def \calR {{\cal R}}
\def \rrb {\overline{\cal R}|}
\def \txtand {{\textrm{~and~}}}
\def \snr {{\textrm{SNR}}} %signal to noise ration
\def \snrdb {{\textrm{SNR}_{\textrm{dB}}}}

%\def \ae {\varepsilon_x}
\def \ae {{E_c}}
\def \ab {\overline{b}}
\def \no { {\cal N}_0 }  %should not it be \N0 with zero
\def \dfreq {e^{j \omega}}

%margin:
%\setlength{\parindent}{0em} \setlength{\oddsidemargin}{0cm}
%\setlength{\evensidemargin}{0cm} \setlength{\textwidth}{16.5cm}

\newcommand{\IR}{\rm I\hspace{-0.06cm}R}
\newcommand{\IN}{\rm I\hspace{-0.06cm}N}

%\newcommand{\ra}{\;\rightarrow\;}
\newcommand{\lra}{\;\Leftrightarrow\;}
\newcommand{\fa}{\;\forall\;}
\newcommand{\ex}{\;\exists\;}

\newcommand{\sinc}{{\textrm{sinc}}}
\newcommand{\rect}{{\textrm{rect}}}
%\def \rect {{\textrm{rect}}} %rect function

\newcommand{\angledeg}[1]{\angle{{#1}}^{\circ}}
\newcommand{\degree}[1]{{{#1}}^{\circ}}